\chapter{Zusammenfassung und Ausblick}

\section{Zusammenfassung}

Die Zusammenfassung fasst die gesamte Arbeit auf typischerweise ein bis zwei Seiten kompakt zusammen. Sie ist neben der Einleitung der meistgelesene Teil — viele Leser lesen nur diese beiden Kapitel. Deshalb sollte die Zusammenfassung eigenständig verständlich sein, ohne dass der Leser andere Kapitel gelesen haben muss.

Eine gute Zusammenfassung folgt der Struktur der Arbeit: Ausgangssituation und Ziel, gewähltes Vorgehen, zentrale Ergebnisse und die wichtigste Schlussfolgerung. Neue Informationen, die nicht in der Arbeit vorkommen, gehören hier nicht hin.

\section{Ausblick}

Der Ausblick beschreibt, welche weiterführenden Schritte nach Abschluss der Arbeit sinnvoll wären. Mögliche Inhalte sind: Welche Funktionen könnten ergänzt werden? Welche zusätzlichen Versuche wären nötig, um die Ergebnisse zu validieren? In welche Richtung könnte die Forschung weiterentwickelt werden?

Der Ausblick zeigt dem Prüfer, dass der Verfasser über den Tellerrand der eigenen Arbeit hinausdenkt.